\documentclass[11pt,a4paper]{article}
\usepackage[utf8]{inputenc}
\usepackage[T1]{fontenc}
\usepackage{geometry}
\usepackage{xcolor}
\usepackage{hyperref}
\usepackage{titlesec}

\geometry{margin=1in}
\definecolor{mtunavy}{RGB}{0,32,91}

\hypersetup{
    colorlinks=true,
    linkcolor=mtunavy,
    urlcolor=mtunavy,
    pdftitle={Cover Letter - Ismail AISSAOUI},
}

\begin{document}

\noindent \textbf{Ismail AISSAOUI} \\
State Engineer in Artificial Intelligence (ESI-SBA) \\
Email: ismail.aissaoui@example.com \\
Phone: [Your Phone Number] \\
LinkedIn: [Your Link]

\bigskip

\begin{flushright}
    Dr. Cemalettin Ozturk \\
    \textbf{Munster Technological University (MTU)} \\
    Department of Computer Science \\
    Cork, Ireland \\
    \medskip
    May 6, 2026
\end{flushright}

\bigskip

\noindent \textbf{Subject: Application for PhD Studentship - Optimisation \& Reasoning (LLMs for Supply Chain)}

\bigskip

Dear Dr. Ozturk,

I am writing to formally submit my application for the PhD Studentship on Large Language Models (LLMs) for Supply Chain Resilience Analysis at MTU. As a State Engineer in Artificial Intelligence from ESI-SBA, I am eager to join your research group to investigate how generative AI can automate complex simulation modeling and provide robust decision intelligence under uncertainty.

My interest in this project is driven by the direct synergy between your objectives and my current research on **Digital Twins** at **Sonatrach**. For my graduation thesis, I am developing a high-fidelity surrogate model for gas turbine monitoring where I hybridize physical laws with advanced sequential architectures, including **Mamba-SSM** and **xLSTM**. This work has given me a deep understanding of:
\begin{itemize}
    \item \textbf{Generative Modeling}: I am proficient in the architectural foundations of Transformers and LLMs. I am particularly interested in how these models can be used to synthesize discrete-event simulations from structured network data, effectively lowering the barrier for expert-level resilience analysis.
    \item \textbf{Uncertainty Quantification (UQ)}: Your project emphasizes "formally quantifying the uncertainty via conformal prediction." In my current work, I develop "Validity Envelopes" for Digital Twins to ensure model reliability. I am keen to apply **Conformal Prediction** and surrogate modeling to provide statistically grounded resilience metrics (Time-to-Survive/Recover).
    \item \textbf{Decision Intelligence}: I believe that LLMs can serve as a bridge between complex simulation data and actionable managerial reasoning, a core theme of your "Optimisation \& Reasoning" research.
\end{itemize}

Beyond my technical skills in Python and PyTorch, I possess the analytical rigour required for a structured, cohort-based PhD. I am a self-driven researcher, and I am excited by the prospect of collaborating with industry partners as part of this Irish national centre for data science and AI.

I have attached my academic transcripts, my CV, and a research statement for your review. I would be honoured to discuss my potential contributions to your project during an interview.

Thank you for your time and consideration.

Sincerely,

\bigskip

\noindent Ismail AISSAOUI

\end{document}
